\section{Background and Threat Model}
\label{sec:background}

\subsection{Provenance Graphs}
\label{sec:bg-provenance}

A system provenance graph $G = (V, E)$ captures the causal history of system execution.
Nodes $V = V_p \cup V_f \cup V_s$ represent three entity types: \emph{processes} ($V_p$), \emph{files} ($V_f$), and \emph{network sockets} ($V_s$).
Directed edges $E \subseteq V \times V \times \mathcal{L}$ represent system call interactions labeled by type $\ell \in \mathcal{L}$, where common labels include \texttt{read}, \texttt{write}, \texttt{execute}, \texttt{connect}, and \texttt{mmap}.
For example, a process reading a configuration file produces an edge $(v_f, v_p, \texttt{read})$; a process spawning a child produces $(v_p, v_{p'}, \texttt{execute})$.

Provenance graphs are constructed from kernel-level audit logs (e.g., Linux Audit, ETW, DTrace) and serve as the foundation for provenance-based HIDS.
Their key property is \emph{completeness}: every system call generates an edge, so attack activities inevitably leave traces in the graph.
This completeness is both the strength of provenance-based detection and the challenge for adversarial evasion---an attacker cannot remove edges without root-level access to the audit subsystem, but can \emph{add} edges by executing additional operations.

\subsection{Detection Layers}
\label{sec:bg-detection}

Modern provenance-based HIDS employ two complementary detection mechanisms.

\paragraph{ML-based detection.}
Systems such as Unicorn~\cite{unicorn}, ShadeWatcher~\cite{shadewatcher}, and ThreaTrace~\cite{threatrace} apply machine learning to provenance graphs.
Unicorn computes graph sketches via histogram-based feature extraction and flags anomalous subgraphs.
ShadeWatcher learns GNN embeddings over the graph structure and detects attacks through embedding-space anomaly scores.
ThreaTrace performs graph-level classification using temporal graph neural networks.
Despite architectural differences, all three systems produce a scalar anomaly score and flag graphs exceeding a learned threshold.

\paragraph{Rule-based detection.}
In addition to ML classifiers, practical deployments incorporate statistical checks that detect structural anomalies in the graph.
These include: (i)~\emph{edge-count thresholds} that flag nodes whose degree exceeds $\mu + k\sigma$ (typically $k=3$) relative to the normal baseline; (ii)~\emph{distribution checks} that measure $\DKL(P_{\text{obs}} \| P_{\text{normal}})$ over operation-type frequencies and flag divergences above a threshold; and (iii)~\emph{temporal burst detectors} that identify unusually dense clusters of operations within sliding time windows.

Prior adversarial evasion work on provenance graphs~\cite{provninja} targets only the ML detection layer.
As we show in \S\ref{sec:evaluation}, camouflage that reduces ML anomaly scores can simultaneously trigger rule-based checks by inflating edge counts or distorting operation distributions.

\subsection{Mimicry Attacks on Provenance Graphs}
\label{sec:bg-mimicry}

Mimicry attacks aim to make malicious provenance subgraphs resemble benign ones by injecting additional operations.
ProvNinja~\cite{provninja} introduced this approach for provenance-based HIDS.
Given an attack subgraph $G_A \subset G$, ProvNinja selects a set of camouflage operations (called \emph{gadgets}) to insert into the host system during attack execution.
Each gadget is a benign system operation---e.g., reading a log file, listing a directory, or querying a DNS name---that adds edges to the provenance graph.

ProvNinja ranks candidate gadgets by a \emph{regularity score} $r(b)$ that measures how frequently operation $b$ appears in normal provenance graphs.
The top-$k$ gadgets by regularity are selected and inserted around the attack operations.
The implicit assumption is that high-regularity operations are both safe (they will not disrupt the attack) and stealthy (they will reduce the anomaly score).

This assumption fails on both counts.
First, regularity does not imply safety: a \texttt{kill -9 \textlangle pid\textrangle} command has high regularity in Linux environments but terminates the attack's host process.
Second, regularity-based selection does not account for the \emph{joint distribution} of operations; inserting many high-regularity operations of the same type skews the operation-type distribution and increases the $\DKL$ divergence that rule-based detectors monitor.

Android-domain evasion techniques face analogous challenges but employ domain-specific solutions.
BagAmmo~\cite{bagammo} wraps injected API calls in try-catch blocks to suppress side effects, and EvadeDroid~\cite{evadedroid} isolates added components from the original APK's execution flow.
AdvDroidZero~\cite{advdroiedzero} uses reinforcement learning to select feature-space perturbations.
These isolation strategies are inapplicable to provenance graphs, where every inserted operation executes at the kernel level with real side effects that are recorded by the audit subsystem.

\subsection{Threat Model}
\label{sec:bg-threat}

We consider an attacker who has gained initial access to a target system and seeks to execute a multi-step attack campaign (e.g., privilege escalation, lateral movement, data exfiltration) while evading provenance-based HIDS.

\paragraph{Attacker capabilities.}
The attacker can execute arbitrary user-level operations on the compromised host, including both attack operations and additional camouflage operations.
Through prior reconnaissance, the attacker has obtained an estimate of the target environment's normal operation distribution $P_{\text{normal}}$---for instance, by observing the system during a benign period or by profiling similar deployments.
The attacker can control the timing and ordering of camouflage operations relative to attack steps.

\paragraph{Attacker limitations.}
The attacker has \emph{black-box} access to the detection system: they know that provenance-based HIDS is deployed but do not know the detector's internal architecture, model weights, or decision thresholds.
The attacker cannot modify or disable the audit subsystem and cannot delete or alter existing edges in the provenance graph.
The attacker cannot observe the detector's real-time output during the attack.

\paragraph{Attack goal.}
The attacker's objective is twofold: (i)~\emph{preserve attack functionality}---all attack steps must execute successfully and achieve their intended effect; and (ii)~\emph{evade detection}---the provenance graph, including both attack and camouflage operations, must not be flagged by either ML-based or rule-based detection mechanisms.
We formalize these requirements in \S\ref{sec:design} through the dependency framework and distribution-matching objective.
